\input{backmatter/checklist-rendered}

\section*{二十六章通路复核}

在前述检查表之外，还应能沿新结构连续回答下面这些跨章问题：

\begin{longtable}{L{0.25\linewidth}L{0.65\linewidth}}
\toprule
通路 & 复核要求 \\
\midrule
RTL 到电路 & 能把组合块、时序块和状态机读回门、寄存器与时钟边界，并区分仿真通过、综合正确和时序收敛。 \\\tablerowrule
ISA 到提交 & 能从机器码指出源状态、数据通路、控制信号、流水阶段、冒险处理和精确异常边界。 \\\tablerowrule
Cache 到页表 & 能把虚拟地址拆成页内偏移与页号，把 Cache 地址拆成 tag/index/offset，并说明 TLB miss、Cache miss 和缺页不是同一事件。 \\\tablerowrule
DRAM 到控制器 & 能从 bank、row buffer、refresh、PHY 训练和调度解释一次主存访问为什么具有可变时延。 \\\tablerowrule
文件到介质 & 能把一次 \code{read()} 写成文件层、块层、提交队列、DMA、控制器、介质、完成队列和唤醒路径。 \\\tablerowrule
PCIe 到设备 & 能解释枚举、BAR、MMIO、TLP、DMA、IOMMU、MSI/MSI-X 和 AER 在同一次设备生命周期中的位置。 \\\tablerowrule
主板到启动 & 能按依赖说明供电、时钟、复位、DDR 训练、设备枚举、固件、装载器和操作系统入口。 \\\tablerowrule
整机故障定位 & 能先判断故障属于电气、时序、协议、地址、所有权、完成还是软件配置，再选择观察点。 \\
\bottomrule
\end{longtable}

\subsection*{怎样完成一次整机复核}

复核时不要只对照术语。先任选一条具体事件，例如一次加载、一次 NVMe 读取、一次网卡收帧或一次缺页，写出事件开始时由谁持有地址、数据与缓冲区；随后逐项记录请求何时被接受、数据经过哪些队列与转换、哪个状态机决定下一步，以及结果在什么边界对 CPU 或设备可见。若只能画出模块名称而无法标出接受、完成和提交时刻，说明通路仍有断点。

第二步主动加入一个故障条件：电源裕量不足、时钟未锁定、地址未映射、credit 耗尽、DMA 映射失效或完成项丢失。先预测最靠近故障源的状态位和计数器，再列出上层最终能看到的症状。故障源、第一条硬件证据和软件症状应当处于同一条时间线上，不能用应用报错直接替代底层证据。

最后交换路径方向，从完成端逆向追踪所有权。例如从中断处理程序回到设备完成队列，再回到 DMA 数据、提交描述符和最初的 doorbell。正向与逆向 trace 能在同一组边界会合，才说明这条通路已经形成可以用于设计和排错的完整模型。
